\documentclass[11pt]{article}

% Change "review" to "final" to generate the final (sometimes called camera-ready) version.
% Change to "preprint" to generate a non-anonymous version with page numbers.
\usepackage[review]{acl}

% Standard package includes
\usepackage{times}
\usepackage{latexsym}

% For proper rendering and hyphenation of words containing Latin characters (including in bib files)
\usepackage[T1]{fontenc}
% This assumes your files are encoded as UTF8
\usepackage[utf8]{inputenc}
% Improves layout and typography
\usepackage{microtype}
% Typewriter font
\usepackage{inconsolata}
% Images
\usepackage{graphicx}

% Added for this manuscript
\usepackage{amsmath}
\usepackage{amssymb}
\usepackage{booktabs}
\usepackage{pgfplots}
\pgfplotsset{compat=1.18}
\usepackage[spanish,es-noshorthands]{babel}

\title{Razonamiento latente adaptativo: \textit{halting} probabilístico tipo PonderNet sobre cadenas de pensamiento implícitas}

\author{Author 1 \\
  Affiliation / Address line 1 \\
  \texttt{email@domain} \\\And
  Author 2 \\
  Affiliation / Address line 1 \\
  \texttt{email@domain} \\\And
  Author 3 \\
  Affiliation / Address line 1 \\
  \texttt{email@domain} \\\AND
  Author 4 \\
  Affiliation / Address line 1 \\
  \texttt{email@domain} \\\And
  Author 5 \\
  Affiliation / Address line 1 \\
  \texttt{email@domain} \\}

\begin{document}
\maketitle

\begin{abstract}
Los modelos de razonamiento en \emph{cadena de pensamiento latente} (\textit{latent chain-of-thought}) reemplazan los pasos de razonamiento en lenguaje natural por una secuencia de vectores continuos, pero fijan de antemano el número de pasos $K$ y lo aplican por igual a toda entrada: un problema de suma de dos pasos consume el mismo presupuesto que uno algebraico de seis. Construyendo sobre la línea Coconut $\rightarrow$ CODI $\rightarrow$ SIM-CoT, injertamos el mecanismo de \textit{halting} probabilístico de PonderNet para que cada instancia decida cuántos pasos latentes necesita: una cabeza de parada define una distribución de terminación sobre los prefijos y un regularizador KL hacia un prior geométrico calibra el presupuesto de cómputo. Sobre GSM8k con un backbone GPT-2, evaluando con decodificación \textit{greedy} y \texttt{batch\_size}$=1$, el modelo adaptativo iguala la exactitud del baseline de $K$ fijo ($\approx 40.8\%$) recortando hasta un $32\%$ los pasos latentes, y aprende un fuerte seguimiento de la dificultad (correlación de Spearman entre número de operaciones y pasos usados de $+0.675$). El aporte central es \textbf{adaptividad de cómputo}: caracterizamos la frontera exactitud--cómputo que los trabajos previos, atados a un único $K$, nunca reportaron.
\end{abstract}

\section{Introducción}

El \emph{chain-of-thought} (CoT) explícito mejora el razonamiento de los modelos de lenguaje generando pasos intermedios en texto \citep{wei2022cot}, pero paga un costo: cada paso se decodifica token a token sobre el vocabulario, y la mayoría de esos tokens existen sólo por fluidez. El \emph{razonamiento latente} (o CoT implícito) elimina ese costo realizando el razonamiento en un espacio continuo: en lugar de proyectar el estado oculto a un token y volver a incrustarlo, el modelo realimenta directamente su último estado oculto como la siguiente incrustación de entrada, produciendo una secuencia de ``pensamientos continuos'' que nunca se decodifican a texto \citep{hao2024coconut}. Esta familia de métodos ---Coconut \citep{hao2024coconut}, CODI \citep{shen2025codi} y SIM-CoT \citep{wei2025simcot}--- alcanza la exactitud del CoT explícito a una fracción del costo de inferencia.

Todos comparten, sin embargo, una limitación estructural: el número de pasos latentes $K$ es un \textbf{hiperparámetro fijo}, idéntico para toda entrada. Un problema de suma de dos pasos y un problema algebraico de seis consumen exactamente el mismo presupuesto de razonamiento. Esto desperdicia cómputo en las entradas fáciles y puede quedarse corto en las difíciles. La premisa de que el cómputo debería escalar con la dificultad ---no con el tamaño de la entrada--- es precisamente la que motiva la literatura de \emph{adaptive computation time} \citep{graves2016act} y, en particular, PonderNet \citep{banino2021pondernet}, que aprende una distribución de parada sobre los pasos de un cómputo recurrente.

Este trabajo hace \textbf{adaptativo por instancia} el número de pasos latentes injertando el mecanismo de \textit{halting} de PonderNet sobre la pila SIM-CoT/CODI. Nuestras contribuciones son:

\begin{enumerate}
  \item Un \textbf{port de PonderNet} (cabeza de parada más regularización hacia un prior geométrico) sobre el bucle de pensamientos continuos de Coconut/CODI/SIM-CoT, decodificando una respuesta candidata tras cada prefijo latente (Sección~\ref{sec:metodo}).
  \item Un \textbf{prior geométrico por instancia} que ancla el presupuesto de parada esperado al número de operaciones de cada problema (Sección~\ref{sec:prior}).
  \item Una \textbf{caracterización de la frontera exactitud--cómputo} ---que los trabajos previos, atados a un único $K$, nunca reportaron--- mostrando recortes de hasta el $32\%$ de pasos a exactitud de paridad (Sección~\ref{sec:frontera}).
\end{enumerate}

Enmarcamos el resultado con honestidad: las ganancias de exactitud sobre el baseline de $K$ fijo están dentro del ruido. El resultado robusto y reproducible es la \textbf{adaptividad de cómputo}: a exactitud de paridad, el modelo asigna menos pasos a los problemas fáciles y más a los difíciles.

\section{Trabajo relacionado}

\paragraph{CoT implícito.} \textbf{Coconut} \citep{hao2024coconut} propone razonar en un espacio latente continuo realimentando el último estado oculto como próxima entrada, y observa que un pensamiento continuo puede codificar varias continuaciones alternativas a la vez (una búsqueda tipo BFS). Se entrena con un currículo multietapa que reemplaza progresivamente los pasos en lenguaje por pensamientos continuos, pero sufre \emph{olvido catastrófico} entre etapas y colapsa al escalar el número de latentes. \textbf{CODI} \citep{shen2025codi} elimina el currículo con una \emph{auto-destilación en una sola etapa}: entrena de forma conjunta una tarea maestra (CoT explícito) y una estudiante (pensamientos continuos) y alinea sus estados ocultos en un único \emph{token de destilación} (la posición justo antes de la respuesta), con \textit{stop-gradient} sobre el maestro. Es el primer CoT implícito que iguala al CoT explícito a escala GPT-2. \textbf{SIM-CoT} \citep{wei2025simcot} diagnostica que la inestabilidad al escalar $K$ proviene de la \emph{homogeneización} de los latentes por falta de supervisión por paso, y añade un \emph{decoder auxiliar} que reconstruye el texto de cada paso de razonamiento durante el entrenamiento; el decoder se descarta en inferencia, de modo que no hay costo adicional. Estabiliza el entrenamiento hasta $K=8$. En los tres métodos, no obstante, $K$ permanece \textbf{fijo} en inferencia.

\paragraph{Cómputo adaptativo.} \textbf{ACT} \citep{graves2016act} introduce el cómputo adaptativo en redes recurrentes mediante una unidad de parada, con gradientes sesgados y un comportamiento sensible a hiperparámetros. \textbf{PonderNet} \citep{banino2021pondernet} lo reformula de manera probabilística: una distribución de parada geométrica sobre los pasos, optimizada con una pérdida esperada más una regularización KL hacia un prior geométrico. No fue diseñado para modelos de lenguaje ni para secuencias de tokens latentes, pero su mecanismo de parada aprendido es exactamente la pieza que falta para volver dinámico el número de pasos en la familia del CoT implícito. Nuestro trabajo hace ese injerto.

\paragraph{Contexto.} Usamos GSM8k \citep{cobbe2021gsm8k} con las anotaciones aumentadas de la internalización paso a paso de \citet{deng2024icot} (GSM8k-Aug), y un backbone GPT-2 \citep{radford2019gpt2}.

\section{Método}
\label{sec:metodo}

\subsection{Fondo: CoT latente de $K$ fijo}

Partiendo del prefijo de la pregunta incrustada $U^{(0)} = (e(x_1),\dots,e(x_T))$, un modelo de CoT latente ejecuta $K$ pasos de razonamiento. En cada paso $k = 1,\dots,K$ toma el estado oculto de la última capa en la última posición y lo realimenta como próxima entrada:
\begin{equation}
  z_k = H_\theta\!\left(U^{(k-1)}\right), \qquad
  U^{(k)} = U^{(k-1)} \oplus z_k,
  \label{eq:latentloop}
\end{equation}
donde $\oplus$ concatena en el eje temporal. Tras $K$ pasos, el modelo conmuta a decodificación sobre el vocabulario y produce la respuesta. El punto crucial es que $K$ se fija de antemano y es idéntico para toda entrada.

\subsection{Halting adaptativo}
\label{sec:halting}

Añadimos una única cabeza lineal de parada, $\texttt{halt\_head}: \mathbb{R}^d \rightarrow \mathbb{R}$, cuyo sesgo se inicializa en $-2.0$ (de modo que la probabilidad de parada inicial es $\sigma(-2)\approx 0.12$ y el modelo usa muchos pasos al comienzo del entrenamiento). En cada paso latente $k$ define una \emph{probabilidad de parada condicional}
\begin{equation}
  \lambda_k = \sigma\!\left(\texttt{halt\_head}(h_k)\right),
\end{equation}
la probabilidad de detenerse en el paso $k$ dado que no se detuvo antes. A diferencia del $K$ fijo, decodificamos una respuesta candidata tras \textbf{cada} prefijo $k = 1,\dots,K$, lo que permite construir la \emph{distribución de parada} geométrica de PonderNet:
\begin{equation}
  p_k = \Bigg(\prod_{j<k}\big(1 - \lambda_j\big)\Bigg)\,\lambda_k .
  \label{eq:halting}
\end{equation}
Para que la distribución sume 1 sin fugar masa, imponemos una \emph{frontera absorbente} en $K$ (equivalente a forzar $\lambda_K \rightarrow 1$): la última entrada recibe toda la masa restante $p_K = \prod_{j<K}(1-\lambda_j)$, y luego se renormaliza.

\subsection{Objetivo de entrenamiento}
\label{sec:loss}

El objetivo combina tres linajes. De PonderNet, la \emph{pérdida esperada de respuesta} y la regularización KL; de SIM-CoT, la supervisión por paso; de CODI, la destilación:
\begin{equation}
  \mathcal{L} = \mathcal{L}_{\text{ponder}}
  + \beta\,\mathcal{L}_{\text{step}}
  + \gamma\,\mathrm{KL}_{\text{geom}}
  + \mathcal{L}_{\text{distill}}
  + \mathcal{L}_{\text{ref}} .
  \label{eq:loss}
\end{equation}
El término principal pondera la entropía cruzada de respuesta de cada prefijo por su probabilidad de parada,
\begin{equation}
  \mathcal{L}_{\text{ponder}}
  = \sum_{k=1}^{K} p_k \, \mathcal{L}_{\text{ans}}^{(k)} ,
  \label{eq:lponder}
\end{equation}
de modo que el modelo aprende a hacer buenas las respuestas de los prefijos a los que asigna masa. $\mathcal{L}_{\text{step}}$ es la reconstrucción del decoder auxiliar de SIM-CoT (con $\beta=1.0$; el decoder se descarta en inferencia). $\mathrm{KL}_{\text{geom}}$ es la divergencia KL de la distribución de parada $p_k$ hacia un \emph{prior geométrico truncado} $q$ de media \texttt{geom\_mean}; con $g = 1/\texttt{geom\_mean}$, $q_k \propto (1-g)^{k}\,g$ con la misma frontera absorbente en $K$. Los términos $\mathcal{L}_{\text{distill}}$ (alineación SmoothL1 de estados ocultos maestro--estudiante) y $\mathcal{L}_{\text{ref}}$ (entropía cruzada de la tarea maestra) provienen de CODI. Por defecto $\gamma=0.01$ y $\texttt{geom\_mean}=3.0$. El peso $\gamma$ es la palanca causal que regula cuánto se acorta el cómputo: sube la presión hacia parar temprano.

\subsection{Prior geométrico por instancia}
\label{sec:prior}

Un prior geométrico \emph{global} pide el mismo presupuesto esperado a todos los problemas. Lo reemplazamos por un prior \textbf{por instancia} que ancla el presupuesto al número de operaciones $n_i$ del problema mediante un mapeo afín:
\begin{equation}
  \texttt{geom\_mean}_i = \mathrm{clamp}\big(\alpha\, n_i + \beta,\; \beta,\; K_{\max}\big).
  \label{eq:adaptiveprior}
\end{equation}
La forma afín (y no la identidad $\texttt{geom\_mean}_i = n_i$) es necesaria: en GSM8k-Aug más de la mitad de los ejemplos tienen $n_i \le 1$, que bajo la identidad caerían en un prior degenerado de masa puntual ($g=1$), cuyo KL domina la pérdida de tarea. El desplazamiento $\beta$ levanta a todos los ejemplos de ese piso.

\subsection{Inferencia}
\label{sec:inferencia}

En inferencia acumulamos masa de parada paso a paso y cortamos el bucle latente en cuanto la masa acumulada cruza un umbral (por defecto $0.5$); $K_{\max} = \texttt{max\_latent\_steps}$ es un tope duro. La respuesta se decodifica desde el prefijo en el que la instancia se detuvo. Un detalle de fidelidad: como los pasos latentes se comparten dentro de un batch pero cada ejemplo se detiene en un paso distinto, una evaluación fiel requiere \texttt{batch\_size}$=1$ (o rebanar la KV-cache por fila en el paso de parada de cada ejemplo); de lo contrario un ejemplo que para temprano se decodificaría desde un prefijo más largo del que su cuenta de pasos reporta. Todos los números de este trabajo usan la variante fiel.

\section{Setup experimental}
\label{sec:setup}

\paragraph{Datos y modelo.} Entrenamos sobre GSM8k-Aug \citep{deng2024icot}, con subconjuntos de 15k/100k ejemplos para el barrido principal (y el conjunto completo, $\approx$385k, para el experimento desde cero de la Sección~\ref{sec:fromscratch}). El backbone es GPT-2 (124M). Salvo que se indique, las corridas parten de un \emph{warm-start} del checkpoint CODI de SIM-CoT (más su decoder auxiliar) y afinan con LoRA. Configuración típica: batch efectivo 128, 5 épocas, \textit{learning rate} $2\times10^{-5}$, $K_{\max}=6$ ($12$ en los experimentos de frontera).

\paragraph{Protocolo de evaluación.} Reportamos exactitud \textbf{\textit{greedy}, de una sola pasada, con \texttt{batch\_size}$=1$}, sobre un conjunto de \textbf{validación de 500 ejemplos}. El baseline de $K$ fijo bajo este protocolo es $\mathbf{40.80\%}$ (\textit{greedy}, $K=6$). Las métricas son: exactitud, número promedio de pasos latentes usados, y la correlación de Spearman entre el número de operaciones del problema ($n_{\text{expr}}$) y los pasos usados, como medida de \emph{seguimiento de dificultad}.

\section{Resultados}
\label{sec:resultados}

\subsection{Exactitud frente al baseline}

Las corridas adaptativas se agrupan en la banda $40.8$--$41.2\%$, es decir, a la par del baseline de $K$ fijo ($40.80\%$) y dentro del ruido de muestreo sobre $n=500$. No reclamamos una mejora de exactitud; el valor del método está en el eje de cómputo.

\subsection{La frontera exactitud--cómputo}
\label{sec:frontera}

El resultado central es que el modelo adaptativo traza una \textbf{frontera} exactitud--pasos que el $K$ fijo no puede: variando $\gamma$ (que hace ``morder'' al prior geométrico) y la forma del prior por instancia, movemos monótonamente el punto de operación. La Tabla~\ref{tab:frontera} compara, por umbral de parada, el mejor baseline adaptativo (prior de pendiente $\alpha=1.0$, $\gamma=0.05$) contra dos reconfiguraciones con $\gamma=0.10$: la \textbf{corrida B} ($\alpha=1.0$, $\beta=1.5$) y la \textbf{corrida C} ($\alpha=0.6$, $\beta=1.5$, que recorta la cola de presupuesto en los problemas difíciles).

\begin{table}[t]
  \centering
  \small
  \begin{tabular}{lccc}
    \toprule
    \textbf{Umbral} & \textbf{Base} & \textbf{Corrida B} & \textbf{Corrida C} \\
                    & $\alpha1.0,\gamma.05$ & $\alpha1.0,\gamma.10$ & $\alpha0.6,\gamma.10$ \\
    \midrule
    0.3 & 39.2 @ 3.27 & 40.0 @ 2.55 & 39.6 @ \textbf{2.03} \\
    0.4 & 40.4 @ 3.78 & 40.6 @ 3.10 & 40.2 @ \textbf{2.47} \\
    0.5 & 40.8 @ 4.34 & \textbf{41.0 @ 3.64} & 40.6 @ \textbf{2.93} \\
    0.8 & 41.0 @ 6.80 & 41.2 @ 6.23 & 40.4 @ 4.85 \\
    \bottomrule
  \end{tabular}
  \caption{Frontera exactitud--cómputo. Cada celda es \emph{exactitud (\%) @ pasos latentes promedio}, en validación ($n=500$), \textit{greedy}, \texttt{batch\_size}$=1$. La corrida B domina al baseline en todos los umbrales; la corrida C recorta drásticamente los pasos a exactitud de paridad.}
  \label{tab:frontera}
\end{table}

\begin{figure}[t]
  \centering
  \begin{tikzpicture}
    \begin{axis}[
      width=\columnwidth, height=5.6cm,
      xlabel={Pasos latentes promedio},
      ylabel={Exactitud (\%)},
      xmin=1.6, xmax=7.2, ymin=38.8, ymax=41.6,
      legend pos=south east, legend cell align=left,
      grid=both, grid style={gray!20},
      tick label style={font=\footnotesize},
      label style={font=\footnotesize},
      legend style={font=\footnotesize},
    ]
      \addplot[mark=*,thick] coordinates {(3.27,39.2)(3.78,40.4)(4.34,40.8)(6.80,41.0)};
      \addlegendentry{Base ($\gamma.05$)}
      \addplot[mark=square*,thick,dashed] coordinates {(2.55,40.0)(3.10,40.6)(3.64,41.0)(6.23,41.2)};
      \addlegendentry{Corrida B ($\gamma.10$)}
      \addplot[mark=triangle*,thick,dotted] coordinates {(2.03,39.6)(2.47,40.2)(2.93,40.6)(4.85,40.4)};
      \addlegendentry{Corrida C ($\gamma.10,\alpha0.6$)}
      \addplot[gray,dashed,thick,domain=1.6:7.2] {40.8};
      \addlegendentry{Baseline $K$ fijo}
    \end{axis}
  \end{tikzpicture}
  \caption{Frontera exactitud--cómputo. Subir $\gamma$ desplaza la frontera hacia la izquierda (menos pasos a igual exactitud); bajar la pendiente $\alpha$ del prior por instancia recorta aún más la cola de presupuesto. La corrida C alcanza $40.6\%$ con $2.93$ pasos, un $32\%$ menos que el baseline de $K$ fijo ($40.80\%$ a $6$ pasos).}
  \label{fig:frontera}
\end{figure}

Los puntos de operación recomendados son claros. Para máxima exactitud, la corrida B a umbral $0.5$ da $\mathbf{41.0\%}$ con $\mathbf{3.64}$ pasos ($+0.2$ pp y $-16\%$ pasos frente al baseline). Para mínimo cómputo, la corrida C a umbral $0.5$ da $\mathbf{40.6\%}$ con $\mathbf{2.93}$ pasos ($-0.2$ pp y $\mathbf{-32\%}$ pasos). En otras palabras, con una pérdida de exactitud despreciable el modelo aprende a resolver GSM8k con menos de la mitad del presupuesto latente del baseline fijo. Un $\gamma$ demasiado bajo ($0.05$) no basta para que el prior muerda; subirlo a $0.10$ mueve toda la frontera. La Figura~\ref{fig:frontera} lo visualiza.

\subsection{Seguimiento de dificultad}
\label{sec:dificultad}

Un modelo verdaderamente adaptativo debería gastar más pasos en los problemas más difíciles. Medimos esto con la correlación de Spearman entre el número de operaciones $n_{\text{expr}}$ y los pasos usados. La correlación crece de forma monótona con la sofisticación del método (Tabla~\ref{tab:spearman}), llegando a $+0.675$ con el afinado de scope completo, prior por instancia y $K_{\max}=12$.

\begin{table}[t]
  \centering
  \small
  \begin{tabular}{lc}
    \toprule
    \textbf{Configuración} & \textbf{Spearman} \\
    \midrule
    Prior global ($\gamma.01$) & $+0.456$ \\
    Barrido $\gamma$ ($\gamma.05$, global) & $+0.553$ \\
    Prior por instancia (\textit{LoRA}) & $+0.662$ \\
    Scope completo $+$ prior $+$ $K_{\max}12$ & $\mathbf{+0.675}$ \\
    \bottomrule
  \end{tabular}
  \caption{Seguimiento de dificultad: correlación de Spearman entre número de operaciones y pasos latentes usados (validación, umbral $0.5$). Crece de forma monótona con la sofisticación del método.}
  \label{tab:spearman}
\end{table}

El efecto es interpretable a nivel de grupo de dificultad: en la mejor configuración, los pasos promedio escalan de $1.83$ (problemas de $0$ operaciones) a $2.50$, $2.86$, $4.49$, $5.43$ y $6.41$ (problemas de $5$ o más operaciones). Los problemas fáciles se detienen alrededor de dos pasos y los difíciles alrededor de seis: el modelo asigna cómputo por dificultad, no por un presupuesto uniforme.

\subsection{Qué no funcionó}

Dos direcciones no dieron fruto. El entrenamiento con \emph{$K$ truncado por instancia} ($K_i = n_i$, forzando a empaquetar cada razonamiento en su cuenta natural de pasos) produce una señal paso-vs-dificultad muy marcada, pero no alcanza la exactitud del prior por instancia en nuestras corridas, por lo que no lo adoptamos; lo atribuimos a una convergencia lenta del \textit{scope} completo, a una regresión de tamaño de batch por memoria y a un $K_i = n_i$ demasiado agresivo. Por otro lado, la variante de prior con desplazamiento $\beta=1.0$ es numéricamente degenerada: pone a los ejemplos de $0$ operaciones (un $16\%$ del entrenamiento) sobre un prior de masa puntual y diverge.

\subsection{Hacia una comparación justa desde cero}
\label{sec:fromscratch}

Los experimentos anteriores parten de un warm-start del checkpoint CODI completo (backbone, LoRA, proyección y decoder, producto de un entrenamiento CODI que no realizamos), y luego afinan sólo cinco épocas. Un revisor puede preguntar, con razón, si el $\approx 41\%$ proviene del \emph{método} adaptativo o del checkpoint heredado. Para aislarlo, replicamos la corrida C desde un GPT-2 \emph{vanilla}: se entrena también el decoder auxiliar (un nuevo \textit{scope} \texttt{full\_dec}), sobre el GSM8k-Aug completo, $40$ épocas. Registramos una corrección de receta necesaria: la tasa $3\times10^{-3}$ del trabajo original diverge en \texttt{bf16} (la norma del gradiente explota y aparece \texttt{NaN}); $1\times10^{-3}$ con \textit{warmup} $0.05$ es estable. Al momento de escribir, este experimento es \textbf{trabajo en curso} (transferido a una GPU A100 por el costo del bucle secuencial de $K=12$). El objetivo es reportar tres puntos comparables ---baseline de $K$ fijo, adaptativo con warm-start, y adaptativo desde cero--- de modo que la brecha cuantifique el aporte del warm-start frente al del método.

\section{Discusión}

Interpretamos la parada aprendida como un mecanismo de \emph{early-exit} que evita el régimen de alto-$K$ propenso al colapso. SIM-CoT \citep{wei2025simcot} documenta que los latentes tardíos se homogeneízan y se limitan a ``repetir'' la predicción final una vez alcanzada la respuesta; detenerse cuando la respuesta ya está codificada preserva la diversidad entre latentes y actúa como regularizador del número de pasos, sin tratarlo como un hiperparámetro a ajustar. Que la exactitud quede dentro del ruido, lejos de ser un resultado negativo, es coherente con esta lectura: el método no busca resolver problemas que el $K$ fijo no resuelve, sino resolver los mismos gastando cómputo proporcional a la dificultad. La frontera de la Figura~\ref{fig:frontera} es la formulación operativa de esa idea, y es precisamente lo que los trabajos previos, cada uno anclado a un único $K$, no podían mostrar.

\section*{Limitaciones}

Varias salvedades acotan nuestras conclusiones. Primero, el conjunto de validación tiene $500$ ejemplos, por lo que diferencias de exactitud de fracciones de punto están dentro del ruido de muestreo; nuestras afirmaciones de exactitud son de \emph{paridad}, no de mejora. Segundo, evaluamos un único backbone (GPT-2) y una única tarea (GSM8k); no sabemos cómo escala la adaptividad a modelos mayores u otros dominios de razonamiento. Tercero, la comparación justa desde cero (Sección~\ref{sec:fromscratch}) aún no tiene número final. Por último, nuestro criterio de parada en inferencia es un umbral determinista sobre la masa acumulada, que difiere del muestreo Bernoulli de la formulación original de PonderNet \citep{banino2021pondernet}; no comparamos ambos regímenes.

\section{Conclusión}

Hicimos adaptativo por instancia el número de pasos de razonamiento latente injertando el \textit{halting} probabilístico de PonderNet sobre la pila CODI/SIM-CoT. En validación, el modelo iguala la exactitud del baseline de $K$ fijo ($\approx 40.8\%$ en GSM8k con GPT-2) recortando hasta un $32\%$ los pasos latentes, y aprende un fuerte seguimiento de la dificultad (Spearman $+0.675$). El aporte no es una exactitud mayor sino \textbf{adaptividad de cómputo}: la caracterización de la frontera exactitud--pasos que la familia del CoT implícito, atada a un $K$ fijo, nunca reportó.

% Custom bibliography entries only
\bibliography{custom}

\end{document}
